\documentclass{ximera}

\begin{document}
	\author{Stitz-Zeager}
	\xmtitle{Exercises for Fundamental Trigonometric Identities}{}

\mfpicnumber{1} \opengraphsfile{ExercisesforFundamentalTrigonometricIdentities} % mfpic settings added 



\begin{question}
In Exercises \ref{useidsforvaluesfirst01} - \ref{useidsforvalueslast01}, use the Reciprocal and Quotient Identities (Theorem \ref{recipquotidfull}) along with the Pythagorean Identities (Theorem \ref{pythids}), to find the value of the circular function requested below. (Find the exact value unless otherwise indicated.)

\begin{problem}\label{useidsforvaluesfirst01}
    If  $\sin(\theta) = \frac{\sqrt{5}}{5}$, find $\csc(\theta)$.

\begin{solution}
$\csc(\theta) = \sqrt{5}$
\end{solution}
\end{problem}

\begin{problem}
    If $\sec(\theta) = - 4$,  find $\cos(\theta)$.

$\cos(\theta) = \answer{-\frac{1}{4}}$
\end{problem}

\begin{problem}
    If $\tan(t) = 3$, find $\cot(t)$.
    
\begin{solution}
$\cot(t) = \frac{1}{3}$
\end{solution}
\end{problem}

\begin{problem}
    If $\theta$ is a Quadrant IV angle with $\cos(\theta) = \frac{5}{13}$, find $\sin(\theta)$.

$\sin(\theta) = \answer{-\frac{12}{13}}$
\end{problem}

\begin{problem}
     If $\theta$ is a Quadrant III angle with $\tan(\theta) = 2$, find $\sec(\theta)$.
\end{problem}

\begin{problem}
    If $\frac{\pi}{2} < t < \pi$ with $\cot(t) = -2$, find $\csc(t)$.
\end{problem}

\begin{problem}
    If $\sec(\theta) = 3$ and $\sin(\theta) < 0$, find $\tan(\theta)$.
\end{problem}

\begin{problem}
     If $\sin(\theta) = -\frac{2}{3}$ but $\tan(\theta) > 0$, find $\cos(\theta)$.
\end{problem}

\begin{problem}
    If $0 < t < \frac{\pi}{2}$ and $\sin(t) = 0.42$, find $\cos(t)$, rounded to four decimal places.
\end{problem}

\begin{problem}
     If $\theta$ is Quadrant IV angle with $\sec(\theta) = 1.17$, find $\tan(\theta)$, rounded to four decimal places.
\end{problem}

\begin{problem}\label{useidsforvalueslast01}  
    If $\pi < t < \frac{3\pi}{2}$ with $\cot(t) = 4.2$, find $\csc(t)$, rounded to four decimal places.
\end{problem}
\end{question}

\begin{question}
In Exercises \ref{useidsforvaluesfirst02} - \ref{useidsforvalueslast02}, use the Reciprocal and Quotient Identities (Theorem \ref{recipquotidfull}) along with the Pythagorean Identities (Theorem \ref{pythids}), to find the exact values of the remaining circular functions.  (Compare your methods with how you solved Exercises \ref{findothercircfirst} - \ref{findothercirclast} in Section \ref{TheOtherCircularFunctions}.)

\begin{problem}\label{useidsforvaluesfirst02}
     $\sin(\theta) = \frac{3}{5}$ with $\theta$ in Quadrant II

$\cos(\theta) = $
\begin{multipleChoice}
  \choice{$\frac{4}{5}$}
  \choice[correct]{$-\frac{4}{5}$}
  \choice{$\frac{3}{4}$}
  \choice{$-\frac{3}{4}$}
\end{multipleChoice}

$\tan(\theta) = $
\begin{multipleChoice}
  \choice{$\frac{4}{5}$}
  \choice{$-\frac{4}{5}$}
  \choice{$\frac{3}{4}$}
  \choice[correct]{$-\frac{3}{4}$}
\end{multipleChoice}

$\csc(\theta) = $
\begin{multipleChoice}
  \choice[correct]{$\frac{5}{3}$}
  \choice{$\frac{5}{4}$}
  \choice{$-\frac{5}{4}$}
  \choice{$-\frac{4}{3}$}
\end{multipleChoice}

$\sec(\theta) = $
\begin{multipleChoice}
  \choice{$\frac{5}{3}$}
  \choice{$\frac{5}{4}$}
  \choice[correct]{$-\frac{5}{4}$}
  \choice{$-\frac{4}{3}$}
\end{multipleChoice}

$\cot(\theta) = $
\begin{multipleChoice}
  \choice{$\frac{5}{3}$}
  \choice{$\frac{5}{4}$}
  \choice{$-\frac{5}{4}$}
  \choice[correct]{$-\frac{4}{3}$}
\end{multipleChoice}

\end{problem}

\begin{problem}
    $\tan(\theta) = \frac{12}{5}$ with $\theta$ in Quadrant III

\begin{solution}
 $\sin(\theta) = -\frac{12}{13}, \cos(\theta) = -\frac{5}{13},  \csc(\theta) = -\frac{13}{12}, \sec(\theta) = -\frac{13}{5}, \cot(\theta) = \frac{5}{12}$
\end{solution}
\end{problem}

\begin{problem}
    $\csc(\theta) = \frac{25}{24}$ with $\theta$ in Quadrant I

$\sin(\theta) = \answer{\frac{24}{25}}$

$\cos(\theta) = \answer{\frac{7}{25}}$

$\tan(\theta) = \answer{\frac{24}{7}}$ 

$\sec(\theta) = \answer{\frac{25}{7}}$

$\cot(\theta) = \answer{\frac{7}{24}}$
\end{problem}

\begin{problem}
    $\sec(\theta) = 7$ with $\theta$ in Quadrant IV
\end{problem}

\begin{problem}
    $\csc(\theta) = -\frac{10\sqrt{91}}{91}$ with $\theta$ in Quadrant III
\end{problem}

\begin{problem}
    $\cot(\theta) = -23$ with $\theta$ in Quadrant II
\end{problem}

\begin{problem}
    $\tan(\theta) = -2$ with $\theta$ in Quadrant IV.
\end{problem}

\begin{problem}
    $\sec(\theta) = -4$ with $\theta$ in Quadrant II.
\end{problem}
 
\begin{problem}
    $\cot(\theta) = \sqrt{5}$ with $\theta$ in Quadrant III.
\end{problem}

\begin{problem}
    $\cos(\theta) = \frac{1}{3}$ with $\theta$ in Quadrant I.
\end{problem}

\begin{problem}
    $\cot(t) = 2$ with $0  < t < \frac{\pi}{2}$.
\end{problem}

\begin{problem}
    $\csc(t) = 5$ with $\frac{\pi}{2} < t < \pi$.
\end{problem}

\begin{problem}
    $\tan(t) = \sqrt{10}$ with $\pi < t < \frac{3\pi}{2}$.
\end{problem}

\begin{problem}\label{useidsforvalueslast02}
     $\sec(t) = 2\sqrt{5}$ with $\frac{3\pi}{2} < t < 2\pi$.
\end{problem}
\end{question}

\begin{problem}
    Skippy claims  $\cos(\theta) + \sin(\theta) = 1$ is an identity because when $\theta = 0$, the equation is true.  Is Skippy correct? Explain.

\begin{solution}
No, Skippy is not correct.  In order to be an identity, an equation must hold for \textit{all} applicable angles.  For example,  $\cos(\theta) + \sin(\theta) = 1$ does not hold when $\theta = \pi$.
\end{solution}
\end{problem}

\begin{question}
In Exercises \ref{firstcirciden} - \ref{lastcirciden}, verify the identity.  Assume that all quantities are defined.

\begin{problem}\label{firstcirciden}
    $\cos(\theta) \sec(\theta) = 1$
\end{problem}

\begin{problem}
    $\tan(t)\cos(t) = \sin(t)$
\end{problem}

\begin{problem}
    $\sin(\theta) \csc(\theta) = 1$
\end{problem}

\begin{problem}
    $\tan(t) \cot(t) = 1$
\end{problem}

\begin{problem}
    $\csc(x) \cos(x) = \cot(x)$
\end{problem}

\begin{problem}
    $\frac{\sin(t)}{\cos^{2}(t)} = \sec(t) \tan(t)$
\end{problem}

\begin{problem}
    $\frac{\cos(\theta)}{\sin^{2}(\theta)} = \csc(\theta) \cot(\theta)$
\end{problem}

\begin{problem}
     $\frac{1+ \sin(x)}{\cos(x)} = \sec(x) + \tan(x)$
\end{problem}

\begin{problem}
    $\frac{1 - \cos(\theta)}{\sin(\theta)} = \csc(\theta) - \cot(\theta)$
\end{problem}

\begin{problem}
     $\frac{\cos(t)}{1 - \sin^{2}(t)} = \sec(t)$
\end{problem}

\begin{problem}
    $\frac{\sin(x)}{1 - \cos^{2}(x)} = \csc(x)$
\end{problem}

\begin{problem}
    $\frac{\sec(t)}{1 + \tan^{2}(t)} = \cos(t)$
\end{problem}

\begin{problem}
    $\frac{\csc(\theta)}{1 + \cot^{2}(\theta)} = \sin(\theta)$
\end{problem}

\begin{problem}
    $\frac{\tan(x)}{\sec^{2}(x) - 1} = \cot(x)$
\end{problem}

\begin{problem}
    $\frac{\cot(t)}{\csc^{2}(t) - 1} = \tan(t)$
\end{problem}

\begin{problem}
     $4 \cos^{2}(\theta) + 4 \sin^{2}(\theta) = 4$
\end{problem}

\begin{problem}
    $9 - \cos^{2}(t) - \sin^{2}(t) = 8$
\end{problem}

\begin{problem}
    $\tan^{3}(t) = \tan(t)\sec^{2}(t) - \tan(t)$
\end{problem}

\begin{problem}
    $\sin^{5}(x) = \left(1-\cos^{2}(x)\right)^{2} \sin(x)$
\end{problem}

\begin{problem}
     $\sec^{10}(t) = \left(1 + \tan^{2}(t)\right)^4 \sec^{2}(t)$
\end{problem}

\begin{problem}
    $\cos^{2}(x)\tan^{3}(x) = \tan(x) - \sin(x)\cos(x)$
\end{problem}

\begin{problem}
    $\sec^{4}(t) - \sec^{2}(t) = \tan^{2}(t) + \tan^{4}(t)$
\end{problem}

\begin{problem}
    $\frac{\cos(\theta) + 1}{\cos(\theta) - 1} = \frac{1 + \sec(\theta)}{1 - \sec(\theta)}$
\end{problem}

\begin{problem}
    $\frac{\sin(t) + 1}{\sin(t) - 1} = \frac{1 + \csc(t)}{1 - \csc(t)}$
\end{problem}

\begin{problem}
    $\frac{1 - \cot(x)}{1+ \cot(x)} = \frac{\tan(x) - 1}{\tan(x) + 1}$
\end{problem}

\begin{problem}
    $\frac{1 - \tan(t)}{1+ \tan(t)} = \frac{\cos(t) - \sin(t)}{\cos(t) + \sin(t)}$
\end{problem}

\begin{problem}
     $\tan(\theta) + \cot(\theta) = \sec(\theta)\csc(\theta)$
\end{problem}

\begin{problem}
    $\csc(t) - \sin(t) = \cot(t)\cos(t)$
\end{problem}

\begin{problem}
    $\cos(x) - \sec(x) = -\tan(x)\sin(x)$
\end{problem}

\begin{problem}
    $\cos(x)(\tan(x) + \cot(x)) = \csc(x)$
\end{problem}

\begin{problem}
    $\sin(t)(\tan(t) + \cot(t)) = \sec(t)$ 
\end{problem}

\begin{problem}
    $\frac{1}{1-\cos(\theta)} + \frac{1}{1+\cos(\theta)} = 2\csc^{2}(\theta)$
\end{problem}

\begin{problem}
    $\frac{1}{\sec(t) + 1} + \frac{1}{\sec(t)-1} = 2 \csc(t) \cot(t)$
\end{problem}

\begin{problem}
    $\frac{1}{\csc(x) + 1} + \frac{1}{\csc(x)-1} = 2 \sec(x) \tan(x)$
\end{problem}

\begin{problem}
    $\frac{1}{\csc(t)-\cot(t)} - \frac{1}{\csc(t) + \cot(t)} = 2 \cot(t)$
\end{problem}

\begin{problem}
    $\frac{\cos(\theta)}{1 - \tan(\theta)} + \frac{\sin(\theta)}{1 - \cot(\theta)} = \sin(\theta) + \cos(\theta)$
\end{problem}

\begin{problem}
    $\frac{1}{\sec(t) + \tan(t)} = \sec(t) - \tan(t)$
\end{problem}

\begin{problem}
    $\frac{1}{\sec(x) - \tan(x)} = \sec(x) + \tan(x)$
\end{problem}

\begin{problem}
    $\frac{1}{\csc(t) - \cot(t)} = \csc(t) + \cot(t)$
\end{problem}

\begin{problem}
    $\frac{1}{\csc(\theta) + \cot(\theta)} = \csc(\theta) - \cot(\theta)$
\end{problem}

\begin{problem}
    $\frac{1}{1-\sin(x)} = \sec^{2}(x) + \sec(x) \tan(x)$
\end{problem}

\begin{problem}
    $\frac{1}{1+\sin(t)} = \sec^{2}(t) - \sec(t) \tan(t)$
\end{problem}

\begin{problem}
    $\frac{1}{1-\cos(\theta)} = \csc^{2}(\theta) + \csc(\theta) \cot(\theta)$
\end{problem}

\begin{problem}
    $\frac{1}{1+\cos(x)} = \csc^{2}(x) - \csc(x) \cot(x)$
\end{problem}

\begin{problem}
    $\frac{\cos(t)}{1 + \sin(t)} = \frac{1-\sin(t)}{\cos(t)}$
\end{problem}

\begin{problem}
    $\csc(\theta) - \cot(\theta) = \frac{\sin(\theta)}{1 + \cos(\theta)}$
\end{problem}

\begin{problem}\label{lastcirciden}
    $\frac{1 - \sin(x)}{1 + \sin(x)} = (\sec(x) - \tan(x))^{2}$
\end{problem}
\end{question}

\begin{question}
In Exercises \ref{logcircidenfirst} - \ref{logcircidenlast}, verify the identity.  You may need to consult Sections \ref{AbsoluteValueFunctions} and \ref{PropertiesofLogarithms} for a review of the properties of absolute value and logarithms before proceeding.

\begin{problem}\label{logcircidenfirst}
    $\quad \ln|\sec(x)| = -\ln|\cos(x)|$ 
\end{problem}

\begin{problem}
     $-\ln|\csc(x)| = \ln|\sin(x)|$
\end{problem}

\begin{problem}
    $-\ln|\sec(x) - \tan(x)| = \ln|\sec(x)+\tan(x)|$
\end{problem}

\begin{problem} \label{logcircidenlast}
     $-\ln|\csc(x) + \cot(x)|= \ln|\csc(x) - \cot(x)|$
\end{problem}
\end{question}

\begin{problem}\label{oneminuscosinelimit}

\begin{enumerate}  

\item What indeterminate form is present in the limit $$\lim_{\theta \rightarrow 0} \frac{1 - \cos(\theta)}{\theta}?$$

\begin{solution}
As $\theta \rightarrow 0$,  $\frac{1 - \cos(\theta)}{\theta} \rightarrow \frac{0}{0}$.
\end{solution}

\item  Graph $f(\theta) = \frac{1 - \cos(\theta)}{\theta}$ near $\theta = 0$. What appears to be $$\lim_{\theta \rightarrow 0} \frac{1 - \cos(\theta)}{\theta}?$$

\begin{solution}
The graph of  $f(\theta) = \frac{1 - \cos(\theta)}{\theta}$ approaches $(0,0)$, so   $\ds{\lim_{\theta \rightarrow 0}}$ $\frac{1 - \cos(\theta)}{\theta}$ appears to be $0$.
\end{solution}

\item\label{oneminuscosinerewrite}  Verify the identity:  $\frac{1 - \cos(\theta)}{\theta} = \frac{\sin(\theta)}{1 + \cos(\theta)} \, \frac{\sin(\theta)}{\theta}$.


\item  Use the fact\footnote{See Exercise \ref{sintovertexercise3} in Section \ref{TheOtherCircularFunctions}.} that $$\lim_{\theta \rightarrow 0} \frac{\sin(\theta)}{\theta} = 1$$ along with  part \ref{oneminuscosinerewrite} to help you find $\lim_{\theta \rightarrow 0} \frac{1 - \cos(\theta)}{\theta}$.

\begin{solution}
$\ds{\lim_{\theta \rightarrow 0}}$ $\frac{1 - \cos(\theta)}{\theta}$ $= \ds{\lim_{\theta \rightarrow 0}}$ $\frac{\sin(\theta)}{1 + \cos(\theta)} \, \frac{\sin(\theta)}{\theta} = \left(\frac{0}{1+\cos(0)}\right)(1) = \left(\frac{0}{2}\right)(1) = 0$.
\end{solution}

\end{enumerate}

\end{problem}
\end{document}
